% !TEX program = xelatex
% ============================================================
% SBW 论文：引言 + 相关工作（核心三小节）重写稿（中文工作稿）
% 依据 2026-07-21 文献调研（docs/paper_draft/review_notes/sbw_related_work_map.md）
% 全部引用经 DBLP / OpenAlex / ACL Anthology / 出版方官网至少一源核验；
% arXiv-only 预印本在文内明示，不作优先权依据。
% 与主稿 sparc_paper_zh.tex 同一导言区，可独立编译；定稿后按节替换回主稿。
% ============================================================
\documentclass[letterpaper]{article}
\usepackage[submission]{aaai2026_prism/aaai2026}
\usepackage{times}
\usepackage{helvet}
\usepackage{courier}
\usepackage[hyphens]{url}
\urlstyle{rm}
\def\UrlFont{\rm}

\usepackage{xeCJK}
\setCJKmainfont[
  BoldFont=SimHei,
  ItalicFont=KaiTi,
  BoldItalicFont=KaiTi
]{SimSun}
\setCJKsansfont[BoldFont=SimHei]{SimHei}
\setCJKmonofont{SimSun}

\usepackage{amsmath,amssymb,amsthm}
\usepackage{booktabs}
\frenchspacing
\setlength{\pdfpagewidth}{8.5in}
\setlength{\pdfpageheight}{11in}
\pdfinfo{/TemplateVersion (2026.1)}

\newcommand{\method}{SPARC}
\newcommand{\sbw}{SBW}
\newcommand{\Sol}{\mathrm{Sol}}
\newcommand{\AnsSol}{\mathrm{Sol}_{A}}
\renewcommand{\refname}{参考文献}

\title{Satisfiable but Wrong: Solver-Decidable Abstention for LLM Constraint Formalization\\
\normalsize（引言与相关工作重写稿）}
\author{匿名作者}

\begin{document}
\maketitle

% ============================================================
\section{引言}
\label{sec:intro}

自然语言约束推理要求系统从题面恢复变量、关系与全局约束，再交由形式求解器
生成答案。Logic-LM~\cite{pan2023}、LINC~\cite{olausson2023} 与
SatLM~\cite{ye2023} 分别将 LLM 与逻辑程序、一阶逻辑证明器和可满足性求解器
结合，这一范式随后扩展到答案集编程~\cite{ishay2023}、约束建模
~\cite{tsouros2023,michailidis2024} 与优化建模~\cite{ahmaditeshnizi2024}。
在基于可满足性求解的流水线中，系统通常将求解器返回 SAT 作为接受答案的依据
~\cite{pan2023,ye2023}。

这一接受规则存在明确盲区：SAT 只说明生成约束彼此一致，不说明它们忠实表达
了题意。当线索被遗漏或弱化时，约束集仍然可满足，而求解器给出的答案在任务
层面是错误的；本文称这类失败为\emph{可满足但错误}
（satisfiable-but-wrong，\sbw{}）。该盲区并非本文杜撰，而是被该方向的
奠基工作反复点名却未解决：Logic-LM 明确承认``合法的符号表示不必然等于
正确的问题形式化''~\cite{pan2023}；其后续改进指出，由于修正环节``没有任何
证明层面的检查''，语法正确但语义错误的公式会被迭代修复直接接受
~\cite{kirtania2024}；SatLM 在附录中记录了求解器正常返回、答案却错误的
案例类别~\cite{ye2023}；在唯一解逻辑谜题上，LLM 生成的答案集程序同样出现
``能运行但解错''的语义错误，补救手段是人工检查~\cite{ishay2023}。
SATBench 进一步表明，线索遗漏后欠约束的形式化仍会稳定返回
SAT~\cite{wei2025}。

现有三类应对均不覆盖这一失败模式。其一，修复流程的触发条件是显式故障——
语法错误、执行失败或 UNSAT：自反馈复审在推理任务上并不稳定
~\cite{madaan2023,huang2024}，而以不可满足核为线索的诊断与修复
~\cite{liffiton2008,chen2024optichat,singh2026}恰好构成 \sbw{} 的对偶——
社区为``过约束导致冲突''建立了成熟的工具链，为``欠约束保持可满足''没有留下
任何触发器。其二，语义对齐路线直接比较题面与形式表示，包括学习式对齐评分
~\cite{lu2025}、候选间符号等价~\cite{li2024symbolic,ryu2025}、回译一致性
~\cite{ying2024}与经形式验证过滤的多次采样投票~\cite{zhou2024dtv}；这些
方法要么依赖额外模型与文本比较，要么需要多倍采样，且信号是启发式或
半可判定的——即便将等价检查确定性化，其召回率也仅约一半
~\cite{poiroux2025}。其三，选择性预测为``何时拒答''提供了标准框架
~\cite{chow1957,geifman2017,kamath2020}，但面向 LLM 的拒答信号大多读取
自评置信度~\cite{xiong2024}或多次采样一致性~\cite{wang2023,farquhar2024}：
前者被系统性证明过度自信~\cite{xiong2024}，后者只度量输出分布的离散度，
对``多次采样一致地遗漏同一条线索''原理上失明。

本文因此研究一个更窄、但由求解器可判定的问题：当任务规范本身给出答案空间
的结构条件时，能否在 SAT 之后、不读取任何金标答案值的前提下，把该条件用作
选择性作答的拒答依据？在唯一答案任务中，该条件即答案变量投影
$\AnsSol(C)$ 恰含一个元素。系统只需一次增量阻塞子句查询——禁止已得的
答案投影并重解——即可判定是否存在第二个不同的答案投影，而无需枚举全部
模型；确定非唯一即结构性拒答，确定唯一则放行，预算内无法判定的
UNKNOWN 与``结构通过''被明确分开。

该信号的理论边界是单侧的，本文将其作为立论前提而非事后限定：若真实答案
唯一，且错误仅来自未排除真实答案的纯欠约束（遗漏或弱化线索），则答案投影
必然非唯一，唯一性门控必然检出；反之，``唯一但错''的形式化会通过门控。
因此该信号是形式化质量的\emph{可判定必要条件}检查，不是语义正确性证明，
也不定位遗漏了哪条线索。诚实地说，构成该门控的每个部件都有先例：阻塞子句
与增量求解是求解器社区的标准原语~\cite{mcmillan2002,eensorensson2004}，
``求一解后禁解重查''在数学应用题的病题检测中已被逐字使用
~\cite{tian2025vcsearch}，SatLM 也已将求解器端的多解与不可满足信号用于
选择性作答~\cite{ye2023}。本文的贡献不在发明机制，而在把散落的实现细节
提炼为一个有理论刻画、有明确语义边界、可系统评测的一等接口。

本文的主要贡献如下：
\begin{itemize}
  \item \textbf{命名并界定 \sbw{} 的可诊断边界。}在唯一答案、纯欠约束的
        前提下，答案投影必然非唯一，故可被唯一性门控检出；唯一但错的
        SAT 输出会逃逸门控。该单侧保证同时给出方法的能与不能。
  \item \textbf{提出 \method{}：SAT 后结构门控的一等接口。}在一般答案
        变量集合 $V_A$ 上构造投影阻塞子句，以单次增量查询实现无需训练、
        不读取金标答案值、由求解器判定的三值输出
        （结构通过 / 结构性拒答 / UNKNOWN）。
  \item \textbf{量化检出范围与覆盖率代价。}在冻结的 ZebraLogic~\cite{lin2025}
        SAT 状态上报告 \sbw{} 检出率与正确输出误拒率，并在风险--覆盖率
        坐标系下与置信度、多次形式化一致性等启发式拒答基线受控比较。
\end{itemize}

% ============================================================
\section{相关工作}
\label{sec:related}

\subsection{LLM 约束形式化与``SAT 即接受''}
\label{sec:rw-formalization}

自动形式化~\cite{wu2022}催生了让 LLM 只负责翻译、由求解器负责推理的
神经符号路线：Logic-LM~\cite{pan2023} 生成逻辑程序并以求解器报错驱动
自我修正，LINC~\cite{olausson2023} 生成一阶逻辑交给证明器并以多次采样
投票稀释翻译噪声，SatLM~\cite{ye2023} 以声明式提示生成可满足性约束；
该范式进一步覆盖答案集编程~\cite{ishay2023}、SMT~\cite{oh2026}、约束
建模~\cite{tsouros2023,michailidis2024}与优化建模~\cite{ahmaditeshnizi2024}。
这些系统的共同接受规则是：求解器一旦成功返回，答案即被采纳；其忠实性
声明均以``形式化正确''为未验证前提~\cite{pan2023}。

与本文最近的三项工作需要逐一划界。\textbf{SatLM}~\cite{ye2023} 已将
求解器端异常用于选择性作答：其错误分类含 UNSAT 与多可行解导致的
AMBIG，遇到两者即拒答并报告选择性准确率。但该歧义检查是面向单个标量
答案变量的实现细节（其 AR-LSAT 路径改用逐选项蕴含检查，机制不同），
论文未将唯一性上升为任务规范给出的结构谓词，未定义一般 $V_A$ 上的投影
阻塞查询，未区分 UNKNOWN 与结构通过，也未给出欠约束与非唯一之间的理论
联系。\textbf{VCSearch}~\cite{tian2025vcsearch} 在数学应用题上逐字实现
了同一求解器机制——求一解、将其作为约束加入、仍可满足即判多解并拒绝——
但其目标是检测\emph{病题}（题面本身缺条件或矛盾），采用多轮 LLM 参与的
迭代搜索，不涉及``好题的坏形式化''，也没有选择性预测评测。
\textbf{SSV}~\cite{raza2025} 同样在不读金标的前提下用求解器验证信号
换取高选择性精度，但其信号是 LLM 生成具体实例后的可满足性检验，每题
需多次模型调用，验证失败先修复重试；本文的唯一性门控只需一次增量求解器
查询，且两类信号正交、可以叠加。此外，近期工作把答案集数量写入训练
奖励以惩罚欠约束的程序~\cite{schrader2026}，但该信号服务于迭代修复
而非拒答。本文相对上述工作的定位是：把``答案投影唯一性''提炼为任务规范
给出的一等结构谓词，配以单侧保证的理论刻画与系统的风险--覆盖率评测。

\subsection{形式化的验证、修复与语义对齐}
\label{sec:rw-repair}

修复路线的触发条件决定了它们的覆盖范围。自反馈复审依赖模型自评
~\cite{madaan2023,shinn2023}，受控研究表明其在推理任务上并不可靠
~\cite{huang2024,olausson2024}；求解器驱动的修复以显式失败为入口——
语法错误与执行异常触发 Logic-LM 式自我修正~\cite{pan2023,kirtania2024}，
UNSAT 与不可满足核触发约束定位与松弛
~\cite{liffiton2008,chen2024optichat,singh2026,ao2026}。触发条件的普查
结论是：没有任何一条修复路线以``SAT 但答案投影非唯一''为触发器；不可
满足核谱系诊断过约束极为成熟，欠约束侧则没有对应工具。软件规格研究
给出了同构的观察：空洞的后置条件（如恒真式）能通过全部正确性检查，
必须另测其判别力~\cite{endres2024}——太弱的规格对一切都满足，正如
欠约束的形式化保持 SAT。

语义对齐路线绕过求解器状态，直接评估题面与形式表示的一致性，按资源
需求可分为：学习式对齐评分需要额外训练的打分模型~\cite{lu2025}；候选
间符号等价需要 $k$ 个候选与两两定理证明~\cite{li2024symbolic}，或在
翻译期用求解器在候选公式间择优~\cite{ryu2025}；回译一致性需要
反向翻译模型与文本蕴含判断~\cite{ying2024}，且遗漏的约束在回译文本中
往往不可见；经形式验证过滤的自一致投票需要 $k$ 倍采样，且验证的是
形式化解的内部一致性而非唯一性~\cite{zhou2024dtv}；自生成测试的执行
共识则受验证器与被验证对象同源之困~\cite{chen2023codet}。这些信号或是
启发式分数，或是``找到证明才可靠''的半可判定判据——即便将等价检查
做成确定性的，召回率也只有约一半，因为自然语言与形式表示的语义等价
本质上不可判定~\cite{poiroux2025}。本文刻意不检查语义等价，只检查任务
规范蕴含的可判定结构条件，换取的是单次求解器查询即可得出的确定结论；
代价是``唯一但错''超出其检测范围——该情形恰是语义对齐方法的适用域，
两条路线互补而非替代。

\subsection{选择性预测与求解器可判定的拒答信号}
\label{sec:rw-selective}

不确定时拒答的框架始于 Chow~\cite{chow1957}，经选择性分类的风险--覆盖率
理论形式化~\cite{elyaniv2010,geifman2017}，并被引入 QA 的分布外场景
~\cite{kamath2020}。面向 LLM 的拒答信号可按证据强度排成谱系：读取自评
置信度的方法被系统性证明过度自信~\cite{xiong2024}，拒答能力也可训练
内化但不提供实例级证据~\cite{zhang2024rtuning}；采样一致性
~\cite{wang2023}与语义熵~\cite{farquhar2024}度量输出分布的离散度，对
``一致地错''结构性失明；共形预测把启发式分数包装为分布层面的边际保证
~\cite{quach2024}，但保证属于分布而非单个实例。近期综述的分类学显示，
由外部求解器证书提供的确定性拒答信号在现有谱系中几乎缺位~\cite{wen2024}
——已知的占位者正是 SatLM 的 UNSAT/AMBIG 拒答~\cite{ye2023}，本文将
其从附带的异常处理提升为有理论刻画的机制。

本文借用的求解器原语均是标准件：阻塞子句枚举出自符号模型检测
~\cite{mcmillan2002}，增量接口使二次查询近乎免费~\cite{eensorensson2004}，
投影变量上的阻塞子句已有成熟实现~\cite{gebser2009}。唯一性检验亦有
正统血统：``给定一解找另一解''被形式化为 ASP 完全问题，动机正是谜题
唯一解检验~\cite{yato2003}，其复杂度定位（唯一可满足性是 coNP 难的
~\cite{valiant1986}）解释了本文只做``证伪唯一性''的廉价方向、并以
UNKNOWN 承接预算耗尽的设计；程序化谜题生成流水线长期以第二解查询
门控生成的谜题实例~\cite{dekegel2020}——区别在于其被门控对象是约束
可信的谜题，本文门控的是约束本身不可信的 LLM 形式化。值得强调的是，
必须在答案变量上投影而非比较全模型：辅助变量的可互换取值会造成与
答案无关的多解~\cite{freuder1991}。概念上更早的先祖是模型检测中的
空洞性检测——规约通过验证、但通过方式暴露了规约本身有病
~\cite{kupferman2003}；本文把同型的``通过即怀疑''检查移到答案投影上。

最后，约束获取与程序合成揭示了本文机制的另一半来路。交互式约束获取
以可信用户回答成员查询而收敛~\cite{bessiere2013,tsouros2024}，
反例引导的综合以完整规格或输入输出 oracle 裁决候选
~\cite{solar2006,jha2010,jha2017}；其中区分性输入查询——问求解器
``是否存在两个都满足当前约束、行为却不同的对象''——与本文的阻塞子句
查询在机制上同源~\cite{jha2010}。分叉点在于裁决：这些方法找到区分性
对象后可以询问可信 oracle 从而继续收敛，而 LLM 形式化场景中不存在
任何可用的 oracle 接口——题面即规格且无法编译为判定器，用户按设定
不具备形式化能力，LLM 自生成的检查约束与被检对象同源
~\cite{chen2023codet}。因此本文只继承机制的查询一半：检测结构性
欠约束即拒答，不声称获取缺失的约束；LLM 生成候选约束收缩解空间仅是
探索性补全，不构成经验证的修复。

% ============================================================
\begin{thebibliography}{99}

% ===== 神经符号形式化管线（2026-07-21 经 DBLP/ACL Anthology 核验）=====
\bibitem{pan2023}
Pan, L.; Albalak, A.; Wang, X.; and Wang, W.~Y. 2023.
Logic-LM: Empowering Large Language Models with Symbolic Solvers for Faithful Logical Reasoning.
In \textit{Findings of EMNLP 2023}, 3806--3824.

\bibitem{olausson2023}
Olausson, T.~X.; Gu, A.; Lipkin, B.; Zhang, C.; Solar-Lezama, A.; Tenenbaum, J.~B.; and Levy, R. 2023.
LINC: A Neurosymbolic Approach for Logical Reasoning by Combining Language Models with First-Order Logic Provers.
In \textit{Proc. EMNLP 2023}, 5153--5176.

\bibitem{ye2023}
Ye, X.; Chen, Q.; Dillig, I.; and Durrett, G. 2023.
SatLM: Satisfiability-Aided Language Models Using Declarative Prompting.
In \textit{Proc. NeurIPS 2023}, vol.~36.

\bibitem{ishay2023}
Ishay, A.; Yang, Z.; and Lee, J. 2023.
Leveraging Large Language Models to Generate Answer Set Programs.
In \textit{Proc. KR 2023}.

\bibitem{kirtania2024}
Kirtania, S.; Gupta, P.; and Radhakrishna, A. 2024.
LOGIC-LM++: Multi-Step Refinement for Symbolic Formulations.
In \textit{Proc. 2nd Workshop on Natural Language Reasoning and Structured
Explanations (NLRSE@ACL 2024)}, 56--63.

\bibitem{oh2026}
Oh, H.; Stern, S.; Lee, Y.; and Philipose, M. 2026.
Neurosymbolic Language Reasoning as Satisfiability Modulo Theory.
arXiv:2602.18095.

\bibitem{tsouros2023}
Tsouros, D.; Verhaeghe, H.; Kad{\i}o\u{g}lu, S.; and Guns, T. 2023.
Holy Grail 2.0: From Natural Language to Constraint Models.
arXiv:2308.01589.

\bibitem{michailidis2024}
Michailidis, K.; Tsouros, D.; and Guns, T. 2024.
Constraint Modelling with LLMs Using In-Context Learning.
In \textit{Proc. CP 2024}, LIPIcs vol.~307, 20:1--20:27.

\bibitem{ahmaditeshnizi2024}
AhmadiTeshnizi, A.; Gao, W.; and Udell, M. 2024.
OptiMUS: Scalable Optimization Modeling with (MI)LP Solvers and Large
Language Models. In \textit{Proc. ICML 2024}, PMLR 235, 577--596.

\bibitem{wei2025}
Wei, A.; Wu, Y.; Wan, Y.; Suresh, T.; Tan, H.; Zhou, Z.; Koyejo, S.;
Wang, K.; and Aiken, A. 2025.
SATBench: Benchmarking LLMs' Logical Reasoning via Automated Puzzle
Generation from SAT Formulas. In \textit{Proc. EMNLP 2025}, 33832--33849.

% ===== 最近先例（2026-07-21 经 ACL Anthology / IJCAI / AAAI 核验）=====
\bibitem{tian2025vcsearch}
Tian, S.-Y.; Zhou, Z.; Yu, K.-Y.; et al. 2025.
Bridging the Gap Between Well-Defined and Ill-Defined Problems in
Mathematical Reasoning. In \textit{Proc. EMNLP 2025}, 12710--12731.

\bibitem{raza2025}
Raza, M.; and Milic-Frayling, N. 2025.
Instantiation-based Formalization of Logical Reasoning Tasks using
Language Models and Logical Solvers.
In \textit{Proc. IJCAI 2025}.

\bibitem{schrader2026}
Schrader, T.~P.; Lange, L.; Kaminski, T.; Razniewski, S.; and
Friedrich, A. 2026.
A Solver-in-the-Loop Framework for Improving LLMs on Answer Set
Programming for Logic Puzzle Solving. In \textit{Proc. AAAI 2026}.

% ===== 修复与诊断 =====
\bibitem{madaan2023}
Madaan, A.; Tandon, N.; Gupta, P.; et al. 2023.
Self-Refine: Iterative Refinement with Self-Feedback.
In \textit{Proc. NeurIPS 2023}, vol.~36.

\bibitem{shinn2023}
Shinn, N.; Cassano, F.; Labash, B.; Gopinath, A.; and Narasimhan, K. 2023.
Reflexion: Language Agents with Verbal Reinforcement Learning.
In \textit{Proc. NeurIPS 2023}, vol.~36.

\bibitem{huang2024}
Huang, J.; Chen, X.; Mishra, S.; Zheng, H.~S.; Yu, A.~W.; Song, X.; and Zhou, D. 2024.
Large Language Models Cannot Self-Correct Reasoning Yet.
In \textit{Proc. ICLR 2024}.

\bibitem{olausson2024}
Olausson, T.~X.; Inala, J.~P.; Wang, C.; Gao, J.; and Solar-Lezama, A. 2024.
Is Self-Repair a Silver Bullet for Code Generation?
In \textit{Proc. ICLR 2024}.

\bibitem{liffiton2008}
Liffiton, M.~H.; and Sakallah, K.~A. 2008. Algorithms for Computing Minimal
Unsatisfiable Subsets of Constraints. \textit{J. Automated Reasoning} 40(1):1--33.

\bibitem{chen2024optichat}
Chen, H.; Constante-Flores, G.~E.; and Li, C. 2024.
Diagnosing Infeasible Optimization Problems Using Large Language Models.
\textit{INFOR: Information Systems and Operational Research} 62(4).

\bibitem{singh2026}
Singh, V.; Cassel, D.; Weir, N.; Feng, N.; and Bayless, S. 2026.
VERGE: Formal Refinement and Guidance Engine for Verifiable LLM Reasoning.
arXiv:2601.20055.

\bibitem{ao2026}
Ao, R.; Simchi-Levi, D.; and Wang, X. 2026.
OptiRepair: Closed-Loop Diagnosis and Repair of Supply Chain Optimization
Models with LLM Agents. arXiv:2602.19439.

\bibitem{endres2024}
Endres, M.; Fakhoury, S.; Chakraborty, S.; and Lahiri, S.~K. 2024.
Can Large Language Models Transform Natural Language Intent into Formal
Method Postconditions? \textit{Proc. ACM Softw. Eng.} 1(FSE).

% ===== 语义对齐（2026-07-21 经 OpenReview / ACL Anthology / PMLR 核验）=====
\bibitem{wu2022}
Wu, Y.; Jiang, A.~Q.; Li, W.; Rabe, M.~N.; Staats, C.; Jamnik, M.; and
Szegedy, C. 2022.
Autoformalization with Large Language Models.
In \textit{Proc. NeurIPS 2022}, vol.~35.

\bibitem{lu2025}
Lu, J.; Wan, Y.; Huang, Y.; Xiong, J.; Liu, Z.; and Guo, Z. 2025.
FormalAlign: Automated Alignment Evaluation for Autoformalization.
In \textit{Proc. ICLR 2025}.

\bibitem{li2024symbolic}
Li, Z.; Wu, Y.; Li, Z.; Wei, X.; Zhang, X.; Yang, F.; and Ma, X. 2024.
Autoformalize Mathematical Statements by Symbolic Equivalence and
Semantic Consistency. In \textit{Proc. NeurIPS 2024}.

\bibitem{ryu2025}
Ryu, H.; Kim, G.; Lee, H.~S.; and Yang, E. 2025.
Divide and Translate: Compositional First-Order Logic Translation and
Verification for Complex Logical Reasoning.
In \textit{Proc. ICLR 2025}.

\bibitem{ying2024}
Ying, H.; Wu, Z.; Geng, Y.; Yuan, Z.; Lin, D.; and Chen, K. 2024.
Lean Workbook: A Large-Scale Lean Problem Set Formalized from Natural
Language Math Problems. In \textit{Proc. NeurIPS 2024 Datasets and
Benchmarks Track}.

\bibitem{zhou2024dtv}
Zhou, J.~P.; Staats, C.; Li, W.; Szegedy, C.; Weinberger, K.~Q.; and
Wu, Y. 2024.
Don't Trust: Verify --- Grounding LLM Quantitative Reasoning with
Autoformalization. In \textit{Proc. ICLR 2024}.

\bibitem{poiroux2025}
Poiroux, A.; Weiss, G.; Kun\v{c}ak, V.; and Bosselut, A. 2025.
Reliable Evaluation and Benchmarks for Statement Autoformalization.
In \textit{Proc. EMNLP 2025}.

\bibitem{chen2023codet}
Chen, B.; Zhang, F.; Nguyen, A.; Zan, D.; Lin, Z.; Lou, J.-G.; and
Chen, W. 2023.
CodeT: Code Generation with Generated Tests.
In \textit{Proc. ICLR 2023}.

% ===== 选择性预测 / LLM 拒答（2026-07-21 核验）=====
\bibitem{chow1957}
Chow, C.~K. 1957.
An Optimum Character Recognition System Using Decision Functions.
\textit{IRE Trans. Electronic Computers} EC-6(4): 247--254.

\bibitem{elyaniv2010}
El-Yaniv, R.; and Wiener, Y. 2010.
On the Foundations of Noise-free Selective Classification.
\textit{J. Machine Learning Research} 11:1605--1641.

\bibitem{geifman2017}
Geifman, Y.; and El-Yaniv, R. 2017.
Selective Classification for Deep Neural Networks.
In \textit{Proc. NeurIPS 2017}, vol.~30.

\bibitem{kamath2020}
Kamath, A.; Jia, R.; and Liang, P. 2020.
Selective Question Answering under Domain Shift.
In \textit{Proc. ACL 2020}, 5684--5696.

\bibitem{xiong2024}
Xiong, M.; Hu, Z.; Lu, X.; Li, Y.; Fu, J.; He, J.; and Hooi, B. 2024.
Can LLMs Express Their Uncertainty? An Empirical Evaluation of Confidence
Elicitation in LLMs. In \textit{Proc. ICLR 2024}.

\bibitem{zhang2024rtuning}
Zhang, H.; Diao, S.; Lin, Y.; Fung, Y.~R.; Lian, Q.; Wang, X.; Chen, Y.;
Ji, H.; and Zhang, T. 2024.
R-Tuning: Instructing Large Language Models to Say `I Don't Know'.
In \textit{Proc. NAACL 2024}.

\bibitem{wang2023}
Wang, X.; Wei, J.; Schuurmans, D.; Le, Q.; Chi, E.; Narang, S.;
Chowdhery, A.; and Zhou, D. 2023.
Self-Consistency Improves Chain of Thought Reasoning in Language Models.
In \textit{Proc. ICLR 2023}.

\bibitem{farquhar2024}
Farquhar, S.; Kossen, J.; Kuhn, L.; and Gal, Y. 2024.
Detecting Hallucinations in Large Language Models Using Semantic Entropy.
\textit{Nature} 630:625--630.

\bibitem{quach2024}
Quach, V.; Fisch, A.; Schuster, T.; Yala, A.; Sohn, J.~H.;
Jaakkola, T.~S.; and Barzilay, R. 2024.
Conformal Language Modeling. In \textit{Proc. ICLR 2024}.

\bibitem{wen2024}
Wen, B.; Yao, J.; Feng, S.; Xu, C.; Tsvetkov, Y.; Howe, B.; and
Wang, L.~L. 2025.
Know Your Limits: A Survey of Abstention in Large Language Models.
\textit{Trans. ACL} 13:529--556.

% ===== 求解器原语与唯一性（2026-07-21 经 OpenAlex / IEICE / AAAI 核验）=====
\bibitem{mcmillan2002}
McMillan, K.~L. 2002.
Applying SAT Methods in Unbounded Symbolic Model Checking.
In \textit{Proc. CAV 2002}, LNCS vol.~2404, 250--264. Springer.

\bibitem{eensorensson2004}
E\'{e}n, N.; and S\"{o}rensson, N. 2004.
An Extensible SAT-solver.
In \textit{SAT 2003 Selected Revised Papers}, LNCS vol.~2919, 502--518. Springer.

\bibitem{gebser2009}
Gebser, M.; Kaufmann, B.; and Schaub, T. 2009.
Solution Enumeration for Projected Boolean Search Problems.
In \textit{Proc. CPAIOR 2009}, LNCS vol.~5547, 71--86. Springer.

\bibitem{yato2003}
Yato, T.; and Seta, T. 2003.
Complexity and Completeness of Finding Another Solution and Its
Application to Puzzles.
\textit{IEICE Trans. Fundamentals} E86-A(5):1052--1060.

\bibitem{valiant1986}
Valiant, L.~G.; and Vazirani, V.~V. 1986.
NP is as Easy as Detecting Unique Solutions.
\textit{Theoretical Computer Science} 47:85--93.

\bibitem{dekegel2020}
De Kegel, B.; and Haahr, M. 2020.
Procedural Puzzle Generation: A Survey.
\textit{IEEE Trans. Games} 12(1):21--40.

\bibitem{freuder1991}
Freuder, E.~C. 1991.
Eliminating Interchangeable Values in Constraint Satisfaction Problems.
In \textit{Proc. AAAI 1991}, 227--233.

% 卷期页码投稿前经 Springer 页面复核
\bibitem{kupferman2003}
Kupferman, O.; and Vardi, M.~Y. 2003.
Vacuity Detection in Temporal Model Checking.
\textit{Int. J. Software Tools for Technology Transfer} 4(2):224--233.

% ===== 约束获取 / 合成与 oracle 依赖（2026-07-21 核验）=====
\bibitem{bessiere2013}
Bessiere, C.; Coletta, R.; Hebrard, E.; Katsirelos, G.; Lazaar, N.;
Narodytska, N.; Quimper, C.-G.; and Walsh, T. 2013.
Constraint Acquisition via Partial Queries.
In \textit{Proc. IJCAI 2013}.

\bibitem{tsouros2024}
Tsouros, D.; Berden, S.; and Guns, T. 2024.
Learning to Learn in Interactive Constraint Acquisition.
In \textit{Proc. AAAI 2024}, 8154--8162.

\bibitem{solar2006}
Solar-Lezama, A.; Tancau, L.; Bodik, R.; Seshia, S.; and Saraswat, V. 2006.
Combinatorial Sketching for Finite Programs.
In \textit{Proc. ASPLOS 2006}, 404--415.

\bibitem{jha2010}
Jha, S.; Gulwani, S.; Seshia, S.~A.; and Tiwari, A. 2010.
Oracle-Guided Component-Based Program Synthesis.
In \textit{Proc. ICSE 2010}, 215--224.

\bibitem{jha2017}
Jha, S.; and Seshia, S.~A. 2017.
A Theory of Formal Synthesis via Inductive Learning.
\textit{Acta Informatica} 54(7):693--726.

% ===== 基准（2026-07-21 经 PMLR / ACL Anthology 核验，author list 已修正）=====
\bibitem{lin2025}
Lin, B.~Y.; Le~Bras, R.; Richardson, K.; Sabharwal, A.; Poovendran, R.;
Clark, P.; and Choi, Y. 2025.
ZebraLogic: On the Scaling Limits of LLMs for Logical Reasoning.
In \textit{Proc. ICML 2025}, PMLR 267, 37889--37905.

\bibitem{zhong2022}
Zhong, W.; Wang, S.; Tang, D.; Xu, Z.; Guo, D.; Chen, Y.; Wang, J.;
Yin, J.; Zhou, M.; and Duan, N. 2022.
Analytical Reasoning of Text.
In \textit{Findings of NAACL 2022}, 2306--2319.

\end{thebibliography}

\end{document}
